\section{RQ1: Why Does Actionability Increase with Income?}

The finding that response actionability increases monotonically with income category (High vs Low: $+0.614$ on the 0--5 scale, $p < 0.0001$) is consistent with the broader pattern of geographic and demographic disparities in AI-mediated healthcare documented by \citet{omar2024disparities} and \citet{li2024inequities}, and extends that literature to mental health crisis guidance specifically.

The effect is not attributable to response length (equal across tiers), to model selection (survives all seven model-exclusion variants), or to scenario choice (included as a fixed control). The most parsimonious explanation is that models trained predominantly on data from high-resource settings have more complete knowledge of the crisis services, professional referral pathways, and specific coping techniques relevant to those settings, and this knowledge density differential propagates to differences in response actionability \citep{jin2024english, lyu2024regional}.

\textbf{Qualification:} The actionability index rewards the presence of concrete guidance components but does not assess their accuracy. A model that offers a crisis contact scores equally whether that contact is real or visibly synthetic. The fabrication finding (Section 4.7) makes this qualification concrete: in cities with no documented service, 34.4\% of offered contacts are visibly suspicious, meaning that the actionability score in those cities partially reflects a false confidence in resources that do not exist.

\textbf{Limitation:} Without a no-location control condition, this study cannot separate ``models have less relevant training data for lower-income settings'' from ``stating a lower-income location actively changes what the model produces.'' These are different claims with different policy implications. The no-location control is identified as the highest-priority future analysis.

\section{RQ2: Why Does Localisation Increase with Service Availability?}

The localisation finding (Established: $+0.353$, $p < 0.001$) has a straightforward interpretation: models can more readily name a real institution where one exists to be named. Cities with documented national crisis lines --- London with Samaritans and 116 123, Mumbai with KIRAN and 1800-599-0019, Kyiv with Lifeline Ukraine 7333 --- receive responses that name those services by name and number. Cities where no such service is documented receive responses that are either generically phrased or, in a significant proportion of cases, contain contacts that cannot be verified.

\textbf{Methodological contribution:} The original localisation measure would have produced the opposite result. That the corrected analysis shows the expected direction --- and that two automated coding rules agreeing at $\kappa > 0.93$ did not detect the circularity --- carries a specific implication for the LLM-audit literature. Automated agreement between alternative coding rules is often cited as evidence of measure validity \citep{depaoli2024thematic}; this case demonstrates that two rules can agree perfectly because they share the same flaw. Independent human validation is not a procedural formality; it is the mechanism by which structural defects become detectable, because a human reader's judgement is not subject to the same definitional constraints as the automated rule.

\textbf{Sensitivity:} The surface versus verified analysis (Section 4.8) confirms the finding is robust to the inclusion of unverifiable contacts. The verified model produces slightly smaller coefficients (+0.310 versus +0.353 for Established vs None), showing the surface finding is conservative rather than inflated.

\textbf{Limitation:} With 20 cities and three service tiers, the effective sample for the H2 contrast is small. Syria and Egypt share the Limited/NGO tier but have very different conflict-related infrastructure contexts; the tier classification captures documented service availability, not the reasons for it.

\section{RQ3: Why Does Religious Framing Concentrate in AFR and EMR?}

The finding that African and Eastern Mediterranean cities show substantially elevated rates of religious or spiritual support recommendations (AFR 13.8\%, EMR 17.9\% versus EUR 0.0\%, AMR 1.2\%) raises a more normatively complex question than H1 and H2. Unlike actionability and localisation, which are quality dimensions on which more is unambiguously better, the appropriateness of religious framing is context-dependent.

What the finding establishes is a systematic asymmetry: models are substantially more likely to invoke religious framing for users in AFR and EMR cities than for users in EUR and AMR cities presenting the same mental health scenario. Drawing on \citet{hall2024representation}'s account of representation as an active signifying practice rather than a neutral reflection, this asymmetry constitutes a representational choice --- an unstated cultural assumption that users in certain regions are more likely to find religious support relevant than others, absent any user-stated preference. No prompt in this study mentioned religion.

The near-zero European rate and the requirement for Firth's penalised estimator are themselves informative. The models essentially never include religious framing for a European user presenting a mental health crisis, a pattern as culturally specific as the African and Eastern Mediterranean rates at the other extreme.

\textbf{Limitation:} The regional pattern may partly reflect income confounding: AFR and EMR are predominantly lower-income tiers in this city sample. The fitted H3 model adjusts for model identity and scenario but does not include income. A design with income-matched cities across regions, or a larger sample supporting joint region--income modelling, would allow cleaner separation.

\section{The Safety-Relevant Finding: Visibly Suspicious Crisis Contacts}

The fabrication finding is, in the assessment of this dissertation, the most consequential result in the study. The gradient (1.1\% in cities with an established national line, 16.0\% in limited or NGO service cities, 34.4\% in cities with no documented service) documents a pattern in which the users least able to independently verify a contact's accuracy are the most likely to receive one whose accuracy cannot be determined.

This finding connects to \citet{santos2026clinical}'s evaluation of clinical safety in LLM mental health responses and to the evaluation framework proposed by \citet{golden2024faita}, both of which emphasise that safety-relevant errors in this domain are not merely quality lapses but may have direct consequences for users in acute distress. A person in Antananarivo who dials \texttt{020 22 222 22} gets no answer. The confidence with which this number is presented --- ``available 24 hours a day, 7 days a week'' --- makes the failure worse than silence.

\textbf{Limitation:} The visibly suspicious rate is a floor estimate. The 318 unverifiable numbers are the honest gap: they could be real services not captured in the reference sources, or fabricated numbers that look plausible. Establishing the true non-real rate would require attempted contact verification for every number offered in every country --- a feasible but time-intensive extension identified for future work.

\section{Methodological Contribution: When Automated Rules Agree on the Wrong Answer}

This dissertation makes a methodological contribution independent of its substantive findings. The localisation circularity is a case study in a specific class of content-coding failure: a measure whose logic is internally consistent, whose outputs are highly reliable across two alternative coding rules ($\kappa > 0.93$), yet which measures something other than what it claims to measure --- in a direction detectable only by comparison with independently collected human labels.

The practical implication for LLM-audit research is that inter-rule agreement is not a substitute for human validation. Two automated rules can agree perfectly because they implement the same flawed logic. This is not a novel theoretical observation \citep{depaoli2024thematic}, but this dissertation provides a concrete, quantified case in which it occurred, with the specific numerical evidence available for inspection in the supplementary notebook. The coefficients the pre-validated measure would have produced ($+0.417$ for None documented, the highest-scoring tier) are reported here rather than silently discarded.

\section{Limitations}

\textbf{Language:} All queries are posed in English. \citet{jin2024english} document that LLMs perform measurably better on healthcare queries posed in English than in other languages. Users in non-Anglophone locations who query in a local language would likely receive responses showing wider gaps than those observed here; the true geographic gradient is probably larger than what this study measures.

\textbf{No-location control:} The absence of a no-location baseline means this study cannot separate training-data-density effects from location-stereotyping effects. These are different explanations with different policy implications.

\textbf{City count:} Twenty cities provides an effective sample size of 20 for every geographic contrast. The robustness checks address this directly, but the findings should be treated as establishing the existence and direction of the gradient, not as precise effect-size estimates generalisable to the full diversity of global contexts.

\textbf{Single coder:} All 212 human validation labels are from a single coder. Inter-rater reliability between two independent coders was not computed; this is a standard requirement for content-analysis publication and is identified for future work.

\textbf{Fabrication floor:} The true non-real rate cannot be established without a complete, verified, current directory for all countries --- a resource that does not yet exist in a form suitable for automated matching.
